\tocsection{\textbf{8/27/26 Lecture 1: }Measure Theory}

Classical probability is a good system, but it is not sufficient for finding the probability of infinite events, such as the probability that it takes an even number of coin flips before getting your first head. This is because we only have finite additivity, so what if we allow countable additivity?\\

\textbf{Countable Additivity:} If $B = \cup_{i = 1}^{\infty} B_{i}$ and the $B_{i}$ are pairwise disjoint, then say
\[ P(B) = \sum_{i = 1}^{\infty} P(B_{i}) \]

However, we run into problems. It turns out sets that we cannot naturally assign a probability to start to exist. Here is an example.\\

\begin{example}{Example 1.1: Vitali Sets}
    Let our set of outcomes be $S^1$, the circle with radius $1$ and circumference $2\pi$. Suppose there is a probability measure $P$ on $\P(S^1)$ that is rotationally invariant; that is, $P(B) = P(B + x)$. Now, we will construct a set that we cannot define a probability to.\\

    Call $\omega', \omega'' \in S^1$ part of the same family if we can write $\omega' = \omega'' + n$ for some $n \in \Z$. Being part of a family is an equivalence relation, so consider the equivalence classes of this relation, and use the AoC to chose a representative from each family. Let $A \subseteq S^1$ be defined as follows
    \[ A = \{\omega \in S^1 \col \omega \text{ is a representative}\} \]
    Now, what is $P(A)$?\\

    First, note that $(A + n) \cap A = \emptyset$ for $n \in \Z$, and that $\cup_{n \in \Z}(A + n) = S^1$. Because of rotational invariance, we see $P(A) = P(A + n)$, but we see
    \[ 1 = P\left(\bigcup_{n \in \Z}(A + n)\right) = \sum_{n \in \Z} P(A + n) = \sum_{n \in \Z} P(A) \]
    However, since $\Z$ is countable, there is no way of assigning a value to $P(A)$ that allows this to be true. 
\end{example}

Thus, we must restrict our probability measuring functions to specific collections of subsets, which motivates the following definition.


\begin{definition}[\textbf{Definition 1.1: }$\sigma$-algebras / $\sigma$-fields]
    Let $\Omega$ be a set. A \emph{$\sigma$-algebra} or \emph{$\sigma$-field} in $\Omega$ is a collection $\mathcal{A} \subseteq \P(\Omega)$ such that, if $A_{i} \in \mathcal{A}$ for $i = 1, 2, \ldots$, then $A_{i}^{c}$, $\cup_{k} A_{k}$, and $\cap_{k} A_{k} \in \mathcal{A}$. Also, $\Omega \in \mathcal{A}$ and $\emptyset \in \mathcal{A}$.
\end{definition}

\begin{remark}{Fun Fact}
    A fun fact is that the ``$\sigma$'' in $\sigma$-algebra and $\sigma$-field stands for ``countable", so whenever you see $\sigma$ in this context, assume countability appears somewhere. 
\end{remark}

Simple examples of a $\sigma$-algebra is $\mathcal{A} = \{\emptyset, \Omega\}$ and $\mathcal{A} = \P(\Omega)$. Here are some more definitions that we wll soon use.

\begin{definition}[\textbf{Definition 1.2: }Measurable Spaces]
    A \emph{measurable space} is a pair $(\Omega, \A)$ where $\A$ is a $\sigma$-algebra. 
\end{definition}

Now, it can be checked quite easily that if $A_{\lambda}$ is a $\sigma$-algebra for each $\lambda \in I$, then $\cap_{\lambda} A_{\lambda}$ is a $\sigma$-algebra. This allows us to make the following definition.

\begin{definition}[\textbf{Definition 1.3: }$\sigma$-algebra Generated by a Set]
    Given a subset $\mathcal{C} \subseteq \P(\Omega)$, $\sigma(\mathcal{C})$ is the intersection of all $\sigma$-algebras containing $\mathcal{C}$. We call $\sigma(\mathcal{C})$ the \emph{$\sigma$-algebra generated by $\mathcal{C}$}. 
\end{definition}

\begin{definition}[\textbf{Definition 1.4: }Borel $\sigma$-algebras]
    Let $S$ be a metric/topological space. If $G$ is the set of all open sets in $S$, let $\mathcal{B}(S) := \sigma(G)$. We call $\mathcal{B}(S)$ the \emph{Borel $\sigma$-algebra} on $S$. 
\end{definition}

\begin{definition}[\textbf{Definition 1.5: }$\pi$-systems]
    A class $\mathcal{C} \subseteq \P(\Omega)$ is called a \emph{$\pi$-system} if it is closed under finite intersections.
\end{definition}

\begin{definition}[\textbf{Definition 1.6: }$\lambda$-systems]
    A class $\mathcal{D} \subseteq \P(\Omega)$ is called a \emph{$\lambda$-system} if it contains $\Omega$ and is closed under proper differences and increasing limits. That is, if $A, B \in \cal{D}$ with $B \subseteq A$, then $A \backslash B \in \cal{D}$, and if $A_1, A_2, \ldots \in \cal{D}$ with $A_1 \uparrow A$ (that is; $A_1 \subseteq A_2 \subseteq \cdots \subseteq A = \cup_{i} A_i$), then $A \in \cal{D}$. 
\end{definition}
